\section{Introduction}\label{sec:intro}

The Black--Litterman model \citep{black1992global} is the standard way to bring investor views into a mean--variance portfolio without the instability of plugging forecasts directly into the optimizer \citep{best1991sensitivity,chopra1993effect}. It needs two inputs per view: the view itself and a statement of how much to trust it, the view uncertainty $\Om$. Practice offers little guidance on the second input. It is usually set by convention, proportional to the asset's variance \citep{he2002intuition}, or elicited as a subjective confidence level \citep{idzorek2007step}.

Large language models (LLMs) produce return forecasts on request \citep{lopez2023can,kim2023if}, and they produce a different forecast each time they are asked. An earlier, unpublished version of this work proposed to use both properties\todo{decide how to refer to the EAAI 2026 submission; it is unpublished}: query the model $N$ times for the same stock and date, take the mean of the answers as the view $q_i$ and their variance as the view uncertainty $\Omega_{ii}$. The appeal is that the forecaster supplies its own confidence, so the input that practice leaves to convention becomes measurable, auditable and scalable across assets. That version reported that Black--Litterman portfolios built this way outperformed equal-weight and mean--variance benchmarks on a large-cap US universe.

This paper re-examines that framework. We change three things that could have produced the earlier result. First, data: the universe is rebuilt at every rebalancing date from point-in-time index membership and market capitalizations in CRSP, removing survivorship and look-ahead in stock selection. Second, models: we use only open-weight LLMs with documented training cutoffs and test strictly after them. Third, evaluation: the Black--Litterman scale parameter $\tau$ is chosen out of sample by a walk-forward rule, and every alternative uses the same rule.

We also make the framework's claim testable. The mapping is useful only if the dispersion it uses as $\Om$ tells the investor which views to trust. We show that evaluating this requires two corrections that the earlier study and much of the LLM-uncertainty literature omit. The realized return differs from the expected return by noise whose variance usually exceeds the view error, so calibration must be judged against the predictive variance $s_i^2+\sigma_i^2/h$, not $s_i^2$ alone. LLM dispersion is also larger for volatile stocks, so the right benchmark is not a constant $\Om$ but the He--Litterman $\Om$, which already scales with volatility. Before running the out-of-sample collection we fixed two tests: (a) dispersion is associated with view error beyond volatility, and (b) holding the views fixed, the dispersion-based $\Om$ gives a better portfolio than the constant and He--Litterman $\Om$.

In the US pilot (S\&P~500, gpt-oss-20b, 2025) test (a) passes and test (b) fails. Dispersion does carry information about forecast error beyond volatility: the within-date rank correlation after volatility standardization is 0.12, positive on 79\% of dates. Used as $\Om$, however, it lowers the Sharpe ratio relative to both conventions in every specification we try, including level matching, a 10\% weight cap, recalibration, four values of the risk-aversion parameter and four cost levels. The reason is a mismatch of magnitudes. Dispersion ranks views roughly correctly but spreads them over two orders of magnitude, while their actual errors differ by a factor of about four, mostly because of volatility. The mapping therefore concentrates the portfolio on a few views that are not more accurate. Behind this lies a more basic finding: with returns as the only input, the model's views are a transformation of short-term momentum (rank correlation 0.96) and have no predictive value at the two-week horizon. With uninformative views no choice of $\Om$ can add value, and one that concentrates does harm.

The paper makes three contributions. First, it provides a point-in-time, cutoff-clean re-examination of the LLM-to-Black--Litterman framework, and shows that the original performance result does not survive it. Second, it gives an evaluation protocol for any source of Black--Litterman view uncertainty: decompose forecast error into view error and return noise, separate shape from level using the $\tau$--$\Om$ invariance, benchmark against He--Litterman, and choose $\tau$ out of sample. Third, it documents why a dispersion signal that is informative in rank can still be harmful as $\Om$, which bears on the broader use of sampled LLM disagreement as a confidence measure in financial decisions.

Section~\ref{sec:related} reviews related work. Section~\ref{sec:framework} states the framework and the tests, Section~\ref{sec:data} the data and design, Section~\ref{sec:results} the results, and Section~\ref{sec:discussion} their implications and limits.
